\section{Projective duality and the molecule modelling equivalence}
\label{sec:molecule-modelling}

The molecule application of the conjecture rests on a chain of three
equivalent rigidity models for a graph $G$ of minimum degree at least
two, all in $\R^3$:
\[
  \text{bar-joint framework of } G^2
  \;\Longleftrightarrow\;
  \text{molecular (hinge-concurrent body-hinge) framework of } G
  \;\Longleftrightarrow\;
  \text{panel-hinge framework of } G .
\]
The first equivalence is the modelling fact that a molecule --- atoms
joined by covalent bonds --- is a bar-joint framework on the square
graph exactly when it is a body-hinge framework whose hinges at each
atom meet at the atom's centre (Whiteley~\cite{whiteley1999}, in the
form used by Jackson--Jord\'an~\cite{jacksonJordan2008}). The second is
projective duality: in $\R^3$ the projective dual sends points to
planes and lines to lines, carrying a hinge-concurrent body-hinge
framework to a panel-hinge framework, and infinitesimal rigidity is
invariant under it (Crapo--Whiteley~\cite{crapoWhiteley1982}). Composed
with Theorem~5.6 (\cref{thm:theorem-55-6}) and the generic
$3$-dimensional bar-joint rigidity matroid of \cref{sec:bar-joint-3d},
this chain yields the molecule rank formula
$r(G^2) = 3|V| - 6 - \operatorname{def}(\tilde G)$
(KT Corollary~5.7), assembled in the next chapter.

This chapter records the two modelling equivalences; the rank formula
that consumes them is the subject of the final chapter. All the
geometry here is specific to $\R^3$: it is where the projective dual
transposes points and planes and where the molecule model lives (KT
\S1.2, and the reconciliation on p.~671).

\subsection{Hinge-concurrent frameworks and the square graph}
\label{sec:molecule-modelling-defs}

\begin{definition}[Hinge-concurrent (molecular) framework]
  \label{def:hinge-concurrent}
  \uses{def:body-hinge-framework}
  A body-hinge framework $(G,p)$ in $\R^3$ is \emph{hinge-concurrent}
  if for every body $v \in V$ there is a point $c(v) \in \R^3$ through
  which every hinge $p(e)$ with $e$ incident to $v$ passes. A
  hinge-concurrent body-hinge framework is also called a
  \emph{molecular framework}~\cite{jacksonJordan2008,whiteley1999},
  since a molecule (atoms as bodies, bonds as hinges) is modelled by
  one: each atom's bonds emanate from the atom's centre.
\end{definition}

\begin{definition}[Square of a graph]
  \label{def:square-graph}
  The \emph{square} of a graph $G = (V,E)$ is the graph
  $G^2 = (V, E^2)$ with
  \[
    E^2 = E \cup \{\, uv \mid u \ne v,\ \exists w \in V \setminus \{u,v\}
      \text{ with } uw, wv \in E \,\},
  \]
  i.e.\ $G$ together with an edge for every pair of vertices at
  distance two.
\end{definition}

\subsection{Projective invariance and duality}
\label{sec:molecule-modelling-projective}

\begin{theorem}[Projective invariance of infinitesimal rigidity; Crapo--Whiteley]
  \label{thm:projective-invariance}
  \uses{def:isInfinitesimallyRigid, def:panel-hinge-framework}
  Infinitesimal rigidity of a framework in $\R^3$ is invariant under
  projective transformation, and in particular under projective
  duality: a framework is infinitesimally rigid if and only if its
  image under a projective transformation is
  (Crapo--Whiteley~\cite{crapoWhiteley1982}, \S3.6). The statement
  applies both to bar-joint frameworks and to (panel-/body-)hinge
  frameworks, whose rigidity is read off the same projectively-defined
  incidences.
\end{theorem}
\begin{proof}
  A projective transformation acts on the homogeneous coordinates of
  the points and on the Pl\"ucker coordinates of the hinges by a fixed
  invertible linear map, transporting the rigidity matrix by pre- and
  post-composition with isomorphisms; the rank, and hence the
  dimension of the trivial-motion complement, is unchanged. The
  detailed argument is Crapo--Whiteley~\cite{crapoWhiteley1982}, \S3.6.
\end{proof}

\begin{lemma}[Projective dual of a panel-hinge framework]
  \label{lem:panel-hinge-dual-molecular}
  \uses{def:panel-hinge-framework, def:hinge-concurrent}
  In $\R^3$, the projective dual of a nonparallel panel-hinge framework
  $(G,p)$ is a hinge-concurrent body-hinge (molecular) framework
  $(G,p^\ast)$ on the same graph, and this correspondence is a
  bijection: the dual of a molecular framework is a panel-hinge
  framework. Under the dual, the panel (a plane) at each body $v$
  becomes the concurrency point $c(v)$, and each hinge line is carried
  to a hinge line through $c(v)$.
\end{lemma}
\begin{proof}
  Projective duality in $\R^3$ transposes points and planes and fixes
  the set of lines, preserving incidence
  (KT~\cite{katohTanigawa2011}, p.~671). The panel $\Pi(v)$ of a
  panel-hinge framework --- a plane --- dualizes to a point $c(v)$; a
  hinge $p(e)$ lying in the panels $\Pi(u)$ and $\Pi(v)$ dualizes to a
  line through both $c(u)$ and $c(v)$. Every hinge incident to $v$ thus
  passes through $c(v)$, which is the hinge-concurrency condition; the
  construction is involutive.
\end{proof}

\subsection{The two modelling equivalences}
\label{sec:molecule-modelling-equivalences}

\begin{theorem}[Panel-hinge $\Leftrightarrow$ molecular realizability]
  \label{thm:panel-hinge-iff-molecular}
  \uses{thm:projective-invariance, lem:panel-hinge-dual-molecular}
  A graph $G$ can be realized as an infinitesimally rigid panel-hinge
  framework in $\R^3$ if and only if it can be realized as an
  infinitesimally rigid hinge-concurrent body-hinge (molecular)
  framework in $\R^3$.
\end{theorem}
\begin{proof}
  Given a rigid panel-hinge realization, its projective dual
  (\cref{lem:panel-hinge-dual-molecular}) is a molecular framework on
  the same graph, and is infinitesimally rigid by projective invariance
  (\cref{thm:projective-invariance}). The converse applies the dual in
  the other direction; the correspondence is a bijection, so no
  realizations are lost.
\end{proof}

\begin{theorem}[Molecule $\Leftrightarrow$ bar-joint of the square; Whiteley, Jackson--Jord\'an]
  \label{thm:molecular-iff-square-bar-joint}
  \uses{def:hinge-concurrent, def:square-graph, def:framework,
        def:isInfinitesimallyRigid}
  Let $G = (V,E)$ be a graph of minimum degree at least two. Then $G$
  can be realized as an infinitesimally rigid molecular
  (hinge-concurrent body-hinge) framework in $\R^3$ if and only if its
  square $G^2$ can be realized as an infinitesimally rigid bar-joint
  framework in $\R^3$.
\end{theorem}
\begin{proof}
  In the square graph a vertex and its neighbours form a complete
  graph, so the bonds at each atom (body) determine a rigid bar-joint
  sub-framework whose motions are exactly the rotations about the lines
  through the atom's centre --- the hinges of the molecular model. This
  is Whiteley's modelling equivalence~\cite{whiteley1999}, for which
  Jackson--Jord\'an~\cite{jacksonJordan2008} is the citable primary
  source (KT resolve the conjecture that closes the argument). The
  minimum-degree-two hypothesis is what makes each atom's neighbourhood
  span enough directions to pin its body, so that the bar-joint
  infinitesimal motions of $G^2$ match the body-hinge motions of $G$
  atom by atom.
\end{proof}
